% Paket dan format naskah: satu kolom, gaya APA 7.

% Ukuran kertas: APA 7 memakai US Letter (8,5 x 11 inci), margin 1 inci.
% Ganti ke a4paper kalau jurnal tujuan meminta A4.
\usepackage[letterpaper,margin=1in]{geometry}

% Encoding dan font
\usepackage[utf8]{inputenc}
\usepackage[T1]{fontenc}
\usepackage{mathptmx}      % mirip Times, termasuk rumus
\usepackage{textcomp}      % simbol tambahan
\usepackage[stretch=30]{microtype}

% Untuk Times New Roman asli: pakai XeLaTeX/LuaLaTeX dan pasang fontnya.
% Nonaktifkan mathptmx, lalu aktifkan dua baris ini:
% \usepackage{fontspec}
% \setmainfont{Times New Roman}

% Bahasa
% Bahasa utama ditaruh terakhir.
\usepackage[english,indonesian]{babel}
% Untuk teks Inggris:
%   \begin{otherlanguage}{english} ... \end{otherlanguage}

% Spasi dan paragraf
\usepackage{setspace}
\doublespacing                       % spasi ganda
\raggedbottom % jangan regangkan isi halaman yang pendek
\usepackage{indentfirst}
\setlength{\parindent}{1.27cm}       % inden 0.5 inci
\setlength{\parskip}{0pt}

% Jangan penggal kata; longgarkan jarak agar teks tetap masuk margin.
\hyphenpenalty=10000
\exhyphenpenalty=10000
\setlength{\emergencystretch}{3em}

% Judul bagian, gaya APA 7
\usepackage{titlesec}
\setcounter{secnumdepth}{0}          % tanpa nomor bagian

% Level 1 -> tengah, tebal
\titleformat{\section}
  {\centering\bfseries\fontsize{12}{14}\selectfont}{}{0em}{}
\titlespacing{\section}{0pt}{*3}{*1.5}

% Level 2 -> rata kiri, tebal
\titleformat{\subsection}
  {\bfseries\fontsize{12}{14}\selectfont}{}{0em}{}
\titlespacing{\subsection}{0pt}{*2.5}{*1}

% Level 3 -> rata kiri, tebal-miring
\titleformat{\subsubsection}
  {\bfseries\itshape\fontsize{12}{14}\selectfont}{}{0em}{}
\titlespacing{\subsubsection}{0pt}{*2}{*0.8}

% Level 4 -> inden, tebal, menyatu dengan paragraf, pakai titik
\titleformat{\paragraph}[runin]
  {\bfseries\fontsize{12}{14}\selectfont}{}{0em}{}[.\ ]
\titlespacing{\paragraph}{1.27cm}{*1.5}{1em}

% Level 5 -> seperti level 4, tapi tebal-miring
\titleformat{\subparagraph}[runin]
  {\bfseries\itshape\fontsize{12}{14}\selectfont}{}{0em}{}[.\ ]
\titlespacing{\subparagraph}{1.27cm}{*1.5}{1em}

% Header: judul pendek dan nomor halaman
\usepackage{fancyhdr}
\pagestyle{fancy}
\fancyhf{}
\newcommand{\runninghead}[1]{\fancyhead[L]{\footnotesize\MakeUppercase{#1}}}
\fancyhead[R]{\footnotesize\thepage}
\renewcommand{\headrulewidth}{0pt}
\setlength{\headheight}{14pt}

% Tabel, gambar, dan caption
\usepackage{graphicx}
\usepackage{float}
\floatplacement{table}{H}
\usepackage{booktabs}
\usepackage{array}
\usepackage{multirow}
\usepackage{caption}
\captionsetup{labelfont=bf, labelsep=period, justification=raggedright,
              singlelinecheck=false, font=small}
\captionsetup[table]{position=above}   % judul di atas tabel

\usepackage{etoolbox}
\AtBeginEnvironment{tabular}{\small}
\setlength{\tabcolsep}{4pt}
\let\origresizebox\resizebox
\renewcommand{\resizebox}[3]{%
  \origresizebox{\ifdim\width>#1 #1\else\width\fi}{#2}{#3}}
% Paksa label bahasa Indonesia meski english jadi bahasa utama TeX.
\renewcommand{\tablename}{Tabel}
\renewcommand{\figurename}{Gambar}

% Rumus dan simbol
\usepackage{amsmath}
\usepackage{amssymb}
\usepackage{icomma}      % koma desimal tanpa spasi ekstra di mode matematika ($0,425$)

% Daftar
\usepackage{enumitem}
\setlist{noitemsep, topsep=4pt}

% Sitasi dan daftar pustaka, gaya APA 7
% Perlu biblatex-apa dan biber.
\usepackage{csquotes}
% Pakai label bibliografi Inggris; naskah tetap bahasa Indonesia.
\DeclareQuoteAlias{english}{indonesian}
\usepackage[
  style=apa,
  backend=biber,
  sorting=nyt,
  sortcites=true,
  uniquename=false,
  uniquelist=false
]{biblatex}
\addbibresource{references.bib}
% Petakan locale Indonesian ke modul english-apa (indonesian-apa.lbx tidak ada).
\DeclareLanguageMapping{indonesian}{english-apa}

% Padanan Indonesia pada sitasi APA 7: "et al." menjadi "dkk." dan "and" pada
% sitasi naratif menjadi "dan" ("&" tetap dipakai pada sitasi dalam kurung).
\DefineBibliographyStrings{english}{andothers = {dkk\adddot}, and = {dan}}
\DefineBibliographyStrings{indonesian}{andothers = {dkk\adddot}, and = {dan}}

% Contoh sitasi:
%   \autocite{key}      -> (Penulis, Tahun)
%   \textcite{key}       -> Penulis (Tahun)
%   \parencite{key}       -> (Penulis, Tahun)
%   \autocite[hlm.\ 12]{key} -> (Penulis, Tahun, hlm. 12)
%
% Aturan APA 7: bila nama penulis menjadi bagian kalimat, tulis nama hanya lewat
% \textcite (Penulis (Tahun), dua penulis "A dan B (Tahun)", tiga atau lebih "A dkk. (Tahun)").
% Bila nama tidak disebut, pakai \autocite atau \parencite ((Penulis, Tahun), dua penulis
% "(A & B, Tahun)"). Jangan menulis nama penulis secara manual di samping sitasi.

% Tautan
% URL panjang boleh pindah baris tanpa tanda hubung.
\usepackage{xurl}
\usepackage[hidelinks]{hyperref} % tanpa kotak warna
\usepackage{orcidlink} % Loads the ORCID link package
\urlstyle{same}

% Longgarkan spasi daftar pustaka untuk DOI/URL panjang.
\AtBeginBibliography{%
  \sloppy
  \setlength{\emergencystretch}{3em}%
}

% Warna
\usepackage{xcolor}   % untuk menandai revisi